\documentclass[../../main.tex]{subfiles}% 注意这里的文件路径不能用 ./main.tex ,否则用latexmk编译子文件会报错
\graphicspath{{\subfix{./image/}}{\subfix{../../image/}}} % 指定图片引用路径,后续可以直接使用图片文件名
% 如果图片存放在上级目录,则使用 ../../image/ 作为备用路径

% 例如:
% \begin{figure}[H]
% \centering
% \includegraphics[width=0.3\textwidth]{图.png}
% \caption{}
% \label{figure:图}
% \end{figure}
% 注意:上述\label{}一定要放在\caption{}之后,否则引用图片序号会只会显示??.

\begin{document}

\section{三角函数相关}

\subsection{三角函数}

\begin{theorem}[两角和差公式]\label{theorem:三角函数和差角公式}
\begin{enumerate}[(1)]
\item \(\sin(\alpha \pm \beta) = \sin\alpha\cos\beta \pm \cos\alpha\sin\beta;\)
\item \(\cos(\alpha \pm \beta) = \cos\alpha\cos\beta \mp \sin\alpha\sin\beta;\)
\item \(\tan(\alpha \pm \beta) = \frac{\tan\alpha \pm \tan\beta}{1 \mp \tan\alpha\tan\beta};\)
\item \(\cot(\alpha \pm \beta) = \frac{\cot\alpha\cot\beta \mp 1}{\cot\beta \pm \cot\alpha}.\)
\end{enumerate}
\end{theorem}
\begin{proof}
\begin{enumerate}[(1)]
\item 在平面直角坐标系中考虑单位圆，设角 \(\alpha\) 和角 \(\beta\) 的终边分别与单位圆交于点 \(P_1\) 和 \(P_2\)，则它们的坐标分别为
\begin{align*}
P_1 = (\cos\alpha, \sin\alpha), \quad P_2 = (\cos\beta, \sin\beta).
\end{align*}
根据向量数量积的坐标表示与几何定义（夹角为 \(\alpha - \beta\)），有
\begin{align*}
\cos(\alpha - \beta) &= P_1 \cdot P_2 = \cos\alpha\cos\beta + \sin\alpha\sin\beta.
\end{align*}
将上式中的 \(\beta\) 替换为 \(-\beta\)，并利用三角函数的奇偶性 \(\sin(-\beta) = -\sin\beta\) 和 \(\cos(-\beta) = \cos\beta\)，可得
\begin{align*}
\cos(\alpha + \beta) = \cos\alpha\cos\beta - \sin\alpha\sin\beta.
\end{align*}
将以上两式合并，即得
\begin{align*}
\cos(\alpha \pm \beta) = \cos\alpha\cos\beta \mp \sin\alpha\sin\beta.
\end{align*}

\item 利用诱导公式 \(\sin\theta = \cos\left(\frac{\pi}{2} - \theta\right)\)，可得
\begin{align*}
\sin(\alpha + \beta) &= \cos\left[ \frac{\pi}{2} - (\alpha + \beta) \right] = \cos\left[ \left(\frac{\pi}{2} - \alpha\right) - \beta \right] \\
&= \cos\left(\frac{\pi}{2} - \alpha\right)\cos\beta + \sin\left(\frac{\pi}{2} - \alpha\right)\sin\beta \\
&= \sin\alpha\cos\beta + \cos\alpha\sin\beta.
\end{align*}
同样用 \(-\beta\) 替换 \(\beta\)，得到
\begin{align*}
\sin(\alpha - \beta) = \sin\alpha\cos\beta - \cos\alpha\sin\beta.
\end{align*}
合并为
\begin{align*}
\sin(\alpha \pm \beta) = \sin\alpha\cos\beta \pm \cos\alpha\sin\beta.
\end{align*}

\item 根据正切与正弦、余弦的商数关系 \(\tan\theta = \frac{\sin\theta}{\cos\theta}\)，代入上述公式可得
\begin{align*}
\tan(\alpha + \beta) = \frac{\sin(\alpha + \beta)}{\cos(\alpha + \beta)} = \frac{\sin\alpha\cos\beta + \cos\alpha\sin\beta}{\cos\alpha\cos\beta - \sin\alpha\sin\beta}.
\end{align*}
将分子分母同时除以 \(\cos\alpha\cos\beta\)，得
\begin{align*}
\tan(\alpha + \beta) = \frac{\tan\alpha + \tan\beta}{1 - \tan\alpha\tan\beta}.
\end{align*}
将 \(\beta\) 替换为 \(-\beta\)，由正切函数的奇偶性 \(\tan(-\beta) = -\tan\beta\)，可得
\begin{align*}
\tan(\alpha - \beta) = \frac{\tan\alpha - \tan\beta}{1 + \tan\alpha\tan\beta}.
\end{align*}
结合以上两式，即得
\begin{align*}
\tan(\alpha \pm \beta) = \frac{\tan\alpha \pm \tan\beta}{1 \mp \tan\alpha\tan\beta}.
\end{align*}

\item 根据余切与正切的倒数关系，并结合 \(\tan\) 的和差公式，将分子分母同时除以 \(\tan\alpha\tan\beta\) 进行化简：
\begin{align*}
\cot(\alpha + \beta) = \frac{1}{\tan(\alpha + \beta)} = \frac{1 - \tan\alpha\tan\beta}{\tan\alpha + \tan\beta} = \frac{\frac{1}{\tan\alpha\tan\beta} - 1}{\frac{1}{\tan\beta} + \frac{1}{\tan\alpha}} = \frac{\cot\alpha\cot\beta - 1}{\cot\beta + \cot\alpha}.
\end{align*}
同样地，可得
\begin{align*}
\cot(\alpha - \beta) = \frac{\cot\alpha\cot\beta + 1}{\cot\beta - \cot\alpha}.
\end{align*}
合并后即得到
\begin{align*}
\cot(\alpha \pm \beta) = \frac{\cot\alpha\cot\beta \mp 1}{\cot\beta \pm \cot\alpha}.
\end{align*}
\end{enumerate}
\end{proof}

\begin{theorem}[倍角公式]\label{theorem:三角函数倍角公式}
\begin{enumerate}[(1)]
\item \(\sin 2\theta = 2\sin\theta\cos\theta;\)
\item \(\cos 2\theta = \cos^2\theta - \sin^2\theta = 2\cos^2\theta - 1 = 1 - 2\sin^2\theta;\)
\item \(\tan 2\theta = \frac{2\tan\theta}{1-\tan^2\theta};\)
\item \(\cot 2\theta = \frac{\cot^2\theta - 1}{2\cot\theta};\)
\item \(\sin 3\theta = 3\sin\theta - 4\sin^3\theta;\)
\item \(\cos 3\theta = 4\cos^3\theta - 3\cos\theta;\)
\item \(\tan 3\theta = \frac{3\tan\theta - \tan^3\theta}{1 - 3\tan^2\theta};\)
\item \(\cot 3\theta = \frac{\cot^3\theta - 3\cot\theta}{3\cot^2\theta - 1};\)
\item \(\sin 4\theta = 4\sin\theta\cos\theta(\cos^2\theta - \sin^2\theta);\)
\item \(\cos 4\theta = 8\cos^4\theta - 8\cos^2\theta + 1.\)
\end{enumerate}
\end{theorem}
\begin{proof}
\begin{enumerate}[(1)]
\item 根据两角和的正弦公式 $\sin(\alpha+\beta) = \sin\alpha\cos\beta + \cos\alpha\sin\beta,$ 令 $\alpha = \beta = \theta,$ 可得
\begin{align*}
\sin 2\theta = \sin\theta\cos\theta + \cos\theta\sin\theta = 2\sin\theta\cos\theta.
\end{align*}
\item 根据两角和的余弦公式 $\cos(\alpha+\beta) = \cos\alpha\cos\beta - \sin\alpha\sin\beta,$ 令 $\alpha = \beta = \theta,$ 可得
\begin{align*}
\cos 2\theta = \cos\theta\cos\theta - \sin\theta\sin\theta = \cos^2\theta - \sin^2\theta.
\end{align*}
再结合同角三角函数关系式 $\sin^2\theta+\cos^2\theta=1,$ 亦可推出
\begin{align*}
\cos 2\theta = 2\cos^2\theta - 1 = 1 - 2\sin^2\theta.
\end{align*}
\item 根据两角和的正切公式 $\tan(\alpha+\beta) = \frac{\tan\alpha + \tan\beta}{1-\tan\alpha\tan\beta},$ 令 $\alpha = \beta = \theta,$ 可得
\begin{align*}
\tan 2\theta = \frac{\tan\theta + \tan\theta}{1 - \tan\theta\cdot\tan\theta} = \frac{2\tan\theta}{1-\tan^2\theta}.
\end{align*}
\item 由正切与余切的倒数关系可知:
\begin{align*}
\cot 2\theta = \frac{1}{\tan 2\theta} = \frac{1 - \tan^2\theta}{2\tan\theta}.
\end{align*}
将其分子分母同时乘以 $\cot^2\theta,$ 即可得到
\begin{align*}
\cot 2\theta = \frac{\cot^2\theta - 1}{2\cot\theta}.
\end{align*}
\item 根据两角和公式 \(\sin(\alpha+\beta) = \sin\alpha\cos\beta+\cos\alpha\sin\beta,\) 令 $\alpha=2\theta, \beta=\theta,$ 可得
\begin{align*}
\sin 3\theta &= \sin(2\theta+\theta) = \sin 2\theta\cos\theta + \cos 2\theta\sin\theta \\
&= 2\sin\theta\cos^2\theta + (1-2\sin^2\theta)\sin\theta \\
&= 2\sin\theta(1-\sin^2\theta) + \sin\theta - 2\sin^3\theta \\
&= 2\sin\theta - 2\sin^3\theta + \sin\theta - 2\sin^3\theta = 3\sin\theta - 4\sin^3\theta.
\end{align*}
\item 根据两角和公式 \(\cos(\alpha+\beta) = \cos\alpha\cos\beta - \sin\alpha\sin\beta,\) 令 $\alpha=2\theta, \beta=\theta,$ 可得
\begin{align*}
\cos 3\theta &= \cos(2\theta+\theta) = \cos 2\theta\cos\theta - \sin 2\theta\sin\theta \\
&= (2\cos^2\theta - 1)\cos\theta - 2\sin^2\theta\cos\theta \\
&= 2\cos^3\theta - \cos\theta - 2(1-\cos^2\theta)\cos\theta \\
&= 2\cos^3\theta - \cos\theta - 2\cos\theta + 2\cos^3\theta = 4\cos^3\theta - 3\cos\theta.
\end{align*}
\item 由两角和的正切公式 \(\tan(\alpha+\beta) = \frac{\tan\alpha + \tan\beta}{1-\tan\alpha\tan\beta},\) 令 $\alpha = 2\theta, \beta = \theta,$ 并设 \(t = \tan\theta\), 可得
\begin{align*}
\tan 3\theta &= \tan(2\theta+\theta) = \frac{\tan 2\theta + \tan\theta}{1 - \tan 2\theta\tan\theta} \\
&= \frac{\frac{2t}{1-t^2} + t}{1 - \frac{2t}{1-t^2} \cdot t} = \frac{\frac{2t+t-t^3}{1-t^2}}{\frac{1-t^2-2t^2}{1-t^2}} = \frac{3t - t^3}{1 - 3t^2} = \frac{3\tan\theta - \tan^3\theta}{1 - 3\tan^2\theta}.
\end{align*}
\item 由余切与正切的倒数关系，并在 $\tan 3\theta$ 的表达式中将分子分母同时乘以 $\cot^3\theta,$ 可得
\begin{align*}
\cot 3\theta = \frac{1}{\tan 3\theta} = \frac{1 - 3\tan^2\theta}{3\tan\theta - \tan^3\theta} = \frac{\cot^3\theta - 3\cot\theta}{3\cot^2\theta - 1}.
\end{align*}
\item 利用二倍角公式 $\sin 2\theta = 2\sin\theta\cos\theta$ 和 $\cos 2\theta = \cos^2\theta-\sin^2\theta,$ 迭代可得
\begin{align*}
\sin 4\theta = 2\sin 2\theta\cos 2\theta = 4\sin\theta\cos\theta(\cos^2\theta - \sin^2\theta).
\end{align*}
\item 利用二倍角公式\(\cos 2\theta =2\cos^2\theta - 1\),迭代可得
\begin{align*}
\cos 4\theta &= 2\cos^2 2\theta - 1 = 2(2\cos^2\theta - 1)^2 - 1 \\
&= 2(4\cos^4\theta - 4\cos^2\theta + 1) - 1 = 8\cos^4\theta - 8\cos^2\theta + 1.
\end{align*}
\item 由二倍角公式\(\cos 2\theta = \cos^2\theta - \sin^2\theta\)可得
\begin{align*}
\cot x - 2\cot 2x &= \frac{\cos x}{\sin x} - 2 \cdot \frac{\cos^2 x - \sin^2 x}{2\sin x\cos x} = \frac{\cos x}{\sin x} - \frac{\cos^2 x - \sin^2 x}{\sin x\cos x} \\
&= \frac{\cos^2 x - (\cos^2 x - \sin^2 x)}{\sin x\cos x} = \frac{\sin^2 x}{\sin x\cos x} = \tan x.
\end{align*}
\end{enumerate}
\end{proof}

\begin{theorem}[三角平方差公式]\label{theorem:三角平方差}
$\sin^2 x-\sin^2 y=\sin(x-y)\sin(x+y)=\cos(y-x)\cos(y+x)=\cos^2 y-\cos^2 x.$
\end{theorem}
\begin{proof}
首先,我们有
\begin{align*}
\cos ^2x-\cos ^2y=1-\sin ^2x-\left( 1-\sin ^2y \right) =\sin ^2y-\sin ^2x.
\end{align*}
接着,我们有
\begin{align*}
\sin(x-y)\sin(x+y) &= (\sin x \cos y - \cos x \sin y)(\sin x \cos y + \cos x \sin y) \\
&= \sin^2 x \cos^2 y - \cos^2 x \sin^2 y \\
&= \sin^2 x (1 - \sin^2 y) - (1 - \sin^2 x) \sin^2 y \\
&= \sin^2 x - \sin^2 y;
\end{align*}
\begin{align*}
\cos(y-x)\cos(y+x) &= (\cos x \cos y + \sin x \sin y)(\cos x \cos y - \sin x \sin y) \\
&= \cos^2 x \cos^2 y - \sin^2 x \sin^2 y \\
&= \cos^2 x \cos^2 y - (1 - \cos^2 x)(1 - \cos^2 y) \\
&= \cos^2 x - \cos^2 y.
\end{align*}
故结论得证.
\end{proof}

\begin{theorem}\label{theorem:sinnx和cosnx以及sin^nx和cos^nx的展开式}
\begin{enumerate}
\item 对$\forall n\in \mathbb{N},$有
\begin{align*}
\sin nx&=\sum_{k=0}^{\lfloor \frac{n-1}{2} \rfloor}{\left( -1 \right) ^k\left( \begin{array}{c}
n\\
2k+1\\
\end{array} \right) (\cos x)^{n-2k-1}(\sin x)^{2k+1}},
\\
\cos nx&=\sum_{k=0}^{\lfloor \frac{n}{2} \rfloor}{\left( -1 \right) ^k\left( \begin{array}{c}
n\\
2k\\
\end{array} \right) (\cos x)^{n-2k}(\sin x)^{2k}},
\\
\tan nx&=\frac{\sin nx}{\cos nx}=\frac{\sum\limits_{k=0}^{\lfloor \frac{n-1}{2} \rfloor}{\left( -1 \right) ^k\left( \begin{array}{c}
n\\
2k+1\\
\end{array} \right) \tan ^{2k+1}x}}{\sum\limits_{r=0}^{\lfloor \frac{n}{2} \rfloor}{\left( -1 \right) ^k\left( \begin{array}{c}
n\\
2k\\
\end{array} \right) \tan ^{2k}x}}.
\end{align*}

\item 对$\forall k\in \mathbb{N} _+,$有
\begin{align*}
&\sin ^{2k-1}x=\frac{1}{2^{2k-2}}\sum_{i=0}^{k-1}{\left( -1 \right) ^{k+i-1}\left( \begin{array}{c}
2k-1\\
i\\
\end{array} \right) \sin \left( \left( 2k-2i-1 \right) x \right)},
\\
&\sin ^{2k}x=\frac{1}{2^{2k-1}}\sum_{i=0}^{k-1}{\left( -1 \right) ^{k+i}\left( \begin{array}{c}
2k\\
i\\
\end{array} \right) \cos \left( \left( 2k-2i \right) x \right)}+\frac{1}{2^{2k}}\left( \begin{array}{c}
2k\\
k\\
\end{array} \right) ,
\\
&\cos ^{2k-1}x=\frac{1}{2^{2k-2}}\sum_{i=0}^{k-1}{\left( \begin{array}{c}
2k-1\\
i\\
\end{array} \right) \cos \left( \left( 2k-2i-1 \right) x \right)},
\\
&\cos ^{2k}x=\frac{1}{2^{2k-1}}\sum_{i=0}^{k-1}{\left( \begin{array}{c}
2k\\
i\\
\end{array} \right) \cos \left( \left( 2k-2i \right) x \right)}+\frac{1}{2^{2k}}\left( \begin{array}{c}
2k\\
k\\
\end{array} \right) .
\end{align*}

\end{enumerate}
\end{theorem}
\begin{remark}
上述结论表明:
$\cos ^nx$可以表示为$1,\cos x,\cdots,\cos nx$的线性组合.
\end{remark}
\begin{proof}
具体证明见\href{https://brilliant.org/wiki/expansions-of-certain-trigonometric-functions/}{Expansions of sin(nx) and cos(nx)}.
% \begin{enumerate}

% \item 

% \item 

% \item 

% \item 

% \item 
% \end{enumerate}
\end{proof}

% \begin{theorem}
% (1) $\cos ^{2n+1}x$可以写成$\cos x,\cos 3x,\cdots ,\cos (2n+1)x$的线性组合 ,即$\cos ^{2n+1}x\in L(\cos x,\cos 3x,\cdots ,\cos (2n+1)x)$ ,也即$\cos ^{2n+1}x=\sum_{k=0}^{n}{a_k\cos (2k+1)x}$ ,其中$a_k\in \mathbb{R}$ ,$k=0,1,\cdots ,n$.

% (2) $\cos ^{2n}x=\frac{1}{2^{2n-1}}\sum_{k=0}^{n-1}{\mathrm{C}_{2n}^{k}\cos 2(n-k)x}+\frac{\mathrm{C}_{2n}^{n}}{2^{2n}}$.

% 综上,$\cos ^nx$可表示为$1,\cos x,\cdots,\cos nx$的线性组合.
% \end{theorem}
% \begin{proof}
% (1) 利用数学归纳法,当$n=1$时,结论显然成立.假设结论对$n-1$成立,则  
% \begin{align*}
% \cos ^{2n+1}x&=\cos ^2x\cdot \cos ^{2n-1}x=\frac{1+\cos 2x}{2}\cdot \sum_{k=0}^{n-1}{a_k\cos \left( 2k+1 \right) x}
% \\
% &=\frac{1}{2}\sum_{k=0}^{n-1}{a_k\cos \left( 2k+1 \right) x}+\frac{1}{2}\sum_{k=0}^{n-1}{a_k\cos 2x\cos \left( 2k+1 \right) x}
% \\
% &=\frac{1}{2}\sum_{k=0}^{n-1}{a_k\cos \left( 2k+1 \right) x}+\frac{1}{2}\sum_{k=0}^{n-1}{a_k\left[ \cos \left( 2k+3 \right) x+\cos \left( 2k-1 \right) x \right]}
% \\
% &=\frac{1}{2}\sum_{k=0}^{n-1}{a_k\cos \left( 2k+1 \right) x}+\frac{1}{2}\sum_{k=1}^{n-1}{a_k\left[ \cos \left( 2k+5 \right) x+\cos \left( 2k+1 \right) x \right]}+\frac{1}{2}a_0\left[ \cos 3x+\cos \left( -x \right) \right] 
% \\
% &=\frac{1}{2}\sum_{k=0}^{n-1}{a_k\cos \left( 2k+1 \right) x}+\frac{1}{2}\sum_{k=1}^{n-1}{a_k\left[ \cos \left( 2k+5 \right) x+\cos \left( 2k+1 \right) x \right]}+\frac{1}{2}a_0\left[ \cos 3x+\cos x \right] .
% \end{align*}  
% 故$\cos ^{2n+1}x\in L(\cos x,\cos 3x,\cdots ,\cos (2n+1)x)$  

% (2) 由二项式定理可得  
% $$
% (1+t^2)^{2n}=\sum_{k=0}^{2n}{\mathrm{C}_{2n}^{k}t^{2k}}
% $$  
% 令$t=e^{\mathrm{i}x}$,则  
% \begin{align*}
% (1+e^{2\mathrm{i}x})^{2n}&=\sum_{k=0}^{2n}{\mathrm{C}_{2n}^{k}e^{2\mathrm{i}kx}} \Rightarrow 2^{2n}\left( \frac{e^{-\mathrm{i}x}+e^{\mathrm{i}x}}{2e^{-\mathrm{i}x}} \right) ^{2n}=\sum_{k=0}^{2n}{\mathrm{C}_{2n}^{k}e^{2\mathrm{i}kx}} \Rightarrow 2^{2n}\left( \frac{e^{-\mathrm{i}x}+e^{\mathrm{i}x}}{2} \right) ^{2n}=e^{-2\mathrm{i}nx}\sum_{k=0}^{2n}{\mathrm{C}_{2n}^{k}e^{2\mathrm{i}kx}}\\
% &\Rightarrow 2^{2n}\cos ^{2n}x=\sum_{k=0}^{2n}{\mathrm{C}_{2n}^{k}e^{2\mathrm{i}(k-n)x}}=\sum_{k=0}^{n-1}[\mathrm{C}_{2n}^{k}e^{2\mathrm{i}(k-n)x}+\mathrm{C}_{2n}^{2n-k}e^{2\mathrm{i}((2n-k)-n)x}]+\mathrm{C}_{2n}^{n}\\
% &=\sum_{k=0}^{n-1}\mathrm{C}_{2n}^{k}(e^{2\mathrm{i}(k-n)x}+e^{2\mathrm{i}(n-k)x})+\mathrm{C}_{2n}^{n}\\
% &\Rightarrow 2^{2n}\cos ^{2n}x=2\sum_{k=0}^{n-1}\mathrm{C}_{2n}^{k}\left( \frac{e^{2\mathrm{i}(k-n)x}+e^{2\mathrm{i}(n-k)x}}{2} \right)+\mathrm{C}_{2n}^{n}=2\sum_{k=0}^{n-1}\mathrm{C}_{2n}^{k}\cos 2(n-k)x+\mathrm{C}_{2n}^{n}\\
% &\Rightarrow \cos ^{2n}x=\frac{1}{2^{2n-1}}\sum_{k=0}^{n-1}\mathrm{C}_{2n}^{k}\cos 2(n-k)x+\frac{\mathrm{C}_{2n}^{n}}{2^{2n}}
% \end{align*}
% 
% \end{proof}







\subsection{反三角函数}

\begin{theorem}[常用反三角函数性质]\label{theorem:常用反三角函数性质}
\begin{enumerate}
\item $$\arcsin x+\arcsin y=\begin{cases}
\arcsin \left( x\sqrt{1-y^2}+y\sqrt{1-x^2} \right) &,xy<0\text{或}x^2+y^2\leqslant 1\\
\pi -\arcsin \left( x\sqrt{1-y^2}+y\sqrt{1-x^2} \right) &,x>0,y>0,x^2+y^2>1\\
-\pi -\arcsin \left( x\sqrt{1-y^2}+y\sqrt{1-x^2} \right) &,x<0,y<0,x^2+y^2>1\\
\end{cases}.$$

\item $$\arcsin x-\arcsin y=\begin{cases}
\arcsin \left( x\sqrt{1-y^2}-y\sqrt{1-x^2} \right) &,xy\geqslant 0\text{或}x^2+y^2\leqslant 1\\
\pi -\arcsin \left( x\sqrt{1-y^2}-y\sqrt{1-x^2} \right) &,x>0,y<0,x^2+y^2>1\\
-\pi -\arcsin \left( x\sqrt{1-y^2}-y\sqrt{1-x^2} \right) &,x<0,y>0,x^2+y^2>1\\
\end{cases}.$$

\item $$\arccos x+\arccos y=\begin{cases}
\arccos \left( xy-\sqrt{1-x^2}\sqrt{1-y^2} \right) &,x+y\geqslant 0\\
2\pi -\arccos \left( xy-\sqrt{1-x^2}\sqrt{1-y^2} \right) &,x+y<0\\
\end{cases}.$$

\item $$\arccos x-\arccos y=\begin{cases}
-\arccos \left( xy+\sqrt{1-x^2}\sqrt{1-y^2} \right) &,x\geqslant y\\
\arccos \left( xy+\sqrt{1-x^2}\sqrt{1-y^2} \right) &,x<y\\
\end{cases}.$$

\item $$\arctan x+\arctan y=\begin{cases}
\arctan \frac{x+y}{1-xy}&,xy<1\\
\pi +\arctan \frac{x+y}{1-xy},x>0&,xy>1\\
-\pi +\arctan \frac{x+y}{1-xy},x<0&,xy>1\\
\end{cases}.$$

\item $$\arctan x-\arctan y=\begin{cases}
\arctan \frac{x-y}{1+xy}&,xy>-1\\
\pi +\arctan \frac{x-y}{1+xy}&,x>0,xy<-1\\
-\pi +\arctan \frac{x-y}{1+xy}&,x<0,xy<-1\\
\end{cases}.$$

\item $$2\arcsin x=\begin{cases}
\arcsin \left( 2x\sqrt{1-x^2} \right) &,\left| x \right|\leqslant \frac{\sqrt{2}}{2}\\
\pi -\arcsin \left( 2x\sqrt{1-x^2} \right) &,\frac{\sqrt{2}}{2}<x\leqslant 1\\
-\pi -\arcsin \left( 2x\sqrt{1-x^2} \right) &,-1\leqslant x<-\frac{\sqrt{2}}{2}\\
\end{cases}.$$

\item $$2\arccos x=\begin{cases}
\arccos \left( 2x^2-1 \right) &,0\leqslant x\leqslant 1\\
2\pi -\arccos \left( 2x^2-1 \right) &,-1\leqslant x<0\\
\end{cases}.$$

\item $$2\arctan x=\begin{cases}
\arctan \frac{2x}{1-x^2},\left| x \right|\leqslant 1\\
\pi +\arctan \frac{2x}{1-x^2}&,\left| x \right|>1\\
-\pi +\arctan \frac{2x}{1-x^2}&,x<-1\\
\end{cases}.$$

\item $$\cos \left( n\arccos x \right) =\frac{\left( x+\sqrt{x^2-1} \right) ^n+\left( x-\sqrt{x^2-1} \right) ^n}{2}\left( n\geqslant 1 \right) .$$
\end{enumerate}
\end{theorem}
\begin{proof}

\end{proof}

\begin{proposition}\label{proposition:arctan相关等式}
$$\arctan x+\arctan \frac{1}{x}=\begin{cases}
\frac{\pi}{2},&x>0\\
-\frac{\pi}{2},&x<0\\
\end{cases}.$$
\end{proposition}
\begin{proof}
令 \( f(x)=\arctan x+\arctan\frac{1}{x} \),则
\begin{align*}
f'(x)&=\frac{1}{x^2 + 1}+\frac{1}{(\frac{1}{x})^2 + 1}(-\frac{1}{x^2})=\frac{1}{x^2 + 1}-\frac{1}{x^2 + 1}=0
\end{align*}
故 \( f(x) \) 为常函数，于是就有 \( f(x)=f(1)=\frac{\pi}{2},\forall x>0 \) ;\( f(x)=f(-1)=-\frac{\pi}{2},\forall x<0 \).
\end{proof}


\subsection{双曲三角函数}

\begin{proposition}
\begin{enumerate}[(1)]
\item $\cosh x=\frac{e^{x}+e^{-x}}{2}\geqslant1,$

\item $\sinh x=\frac{e^{x}-e^{-x}}{2}\geqslant x.$
\end{enumerate}
\end{proposition}
\begin{proof}
可以分别利用均值不等式和求导进行证明. 
\end{proof}

\begin{proposition}\label{proposition:双曲三角函数基本性质}
\begin{enumerate}
\item $\cosh^2 x-\sinh^2 x=1$.

\item $\cosh(2x)=2\cosh^2 x-1=1-2\sinh^2 x=\cosh^2 x+\sinh^2 x$.

\item $\sinh(2x)=2\sinh x\cosh x$. 
\end{enumerate}
\end{proposition}
\begin{proof}

\end{proof}




\end{document}